\documentclass[10pt]{article}
\usepackage[margin=2.2cm]{geometry}
\usepackage{booktabs,longtable,amsmath,amssymb,bm,graphicx}
\usepackage[hidelinks]{hyperref}
\usepackage{times}
% the generated tables use the Bioinformatics template macros
\newcommand{\processtable}[3]{\caption{#1}\centering\small #2 #3}
\newcommand{\botrule}{\bottomrule}
\renewcommand{\thetable}{S\arabic{table}}
\renewcommand{\thefigure}{S\arabic{figure}}
\renewcommand{\thesection}{S\arabic{section}}
\makeatletter\let\tablestar\table\let\endtablestar\endtable
\renewenvironment{table*}{\begin{table}}{\end{table}}\makeatother

\title{Supplementary information for\\``Graph-conditioned state-space models for genetic perturbation
response prediction: a controlled evaluation against simple baselines''}
\author{Aryan Padarthi and Riley Huynh}
\date{}

\begin{document}
\maketitle

\section{Legacy perturbation encoding}
Table~\ref{tab:audit} quantifies the encoding flaw described in Section~3.1 of the main text, computed
with the original data loader (\texttt{legacy/data\_loader.py}) and split manifests
(\texttt{splits/splits\_manifest\_*.json}) by \texttt{gcpm/audit\_hvg\_encoding.py}. HVGs were selected
on training cells only, as in the original pipeline.

\begin{table}[h]
\caption{Legacy HVG-restricted encoding (Section~3.1): number of the 105 CRISPRa targets among the 500 HVGs, and perturbations whose input vector is all zero, for the ten split manifests of the earlier version of this work.}\label{tab:audit}
\centering\small
\begin{tabular}{rrrrrr}
\toprule
Split & Targets in HVGs & Train all-zero & Distinct train inputs & Test all-zero & Test $n$ \\
\midrule
0 & 17 & 133/169 (79\%) & 16 & 52/66 (79\%) & 66 \\
1 & 12 & 146/183 (80\%) & 14 & 49/52 (94\%) & 52 \\
2 & 14 & 147/174 (84\%) & 14 & 43/61 (70\%) & 61 \\
3 & 13 & 133/161 (83\%) & 14 & 60/74 (81\%) & 74 \\
4 & 13 & 134/167 (80\%) & 15 & 61/68 (90\%) & 68 \\
5 & 14 & 132/168 (79\%) & 17 & 59/67 (88\%) & 67 \\
6 & 14 & 133/161 (83\%) & 15 & 58/74 (78\%) & 74 \\
7 & 15 & 136/177 (77\%) & 19 & 50/58 (86\%) & 58 \\
8 & 14 & 137/167 (82\%) & 14 & 58/68 (85\%) & 68 \\
9 & 14 & 135/176 (77\%) & 19 & 55/59 (93\%) & 59 \\
\bottomrule
\end{tabular}
\end{table}


\section{Model development on the validation set of split 1}
All design decisions were made on the validation perturbations of split 1; test perturbations were not
used. Table~\ref{tab:pilot} lists the successive changes and the lowest validation MSE (all genes)
reached in short pilot runs (12--35 epochs). For reference, predicting no change gives 0.00433, the
train mean 0.00361, the linear model 0.00312 and the additive baseline 0.00244 on the same perturbations
(the latter using the training singles only).

\begin{table}[h]
\caption{Pilot runs of GCP-Mamba on the validation perturbations of split 1.}\label{tab:pilot}
\centering\small
\begin{tabular}{lll}
\toprule
Change & Best validation MSE & Observation\\
\midrule
Initial model (scalar readout, raw PCA loadings) & 0.00362 & equal to the train mean; training loss also at train-mean level\\
+ bilinear readout & 0.00360 & unchanged\\
+ standardised gene embeddings & 0.00569 & fits training data, generalises worse than the mean\\
+ per-target embeddings, dropout, weight decay $10^{-2}$ & 0.00523 & large initial loss (random outputs on all genes)\\
+ initialisation at the train-mean prediction ($\text{lr}=2\times10^{-3}$) & 0.00320 & first improvement over the mean\\
+ learning rate $10^{-3}$ (final) & 0.00297 & \\
\bottomrule
\end{tabular}
\end{table}

\section{Statistical comparisons and initialisation variance}
Table~\ref{tab:tests} lists all pairwise tests between GCP-Mamba and the other models on the Norman
data (deep models averaged over initialisations). Table~\ref{tab:reps} shows how much the split-level
error of each model changes between initialisations; the within-split range is of the same size as
the differences between graph-conditioned and graph-free variants. Table~\ref{tab:pearson} gives the
Pearson correlation on the top-20 DE genes for every model.

\begin{longtable}{llrrrr}
\caption{GCP-Mamba versus every other model on the Norman data: two-sided Wilcoxon signed-rank tests on per-perturbation differences in MSE (top-20 DE genes), pooled over the five splits; Holm correction within each test group. Negative median differences favour GCP-Mamba.}\label{tab:tests}\\
\toprule
Test group & Comparison model & $n$ & Median diff. & $P$ & $P_{\mathrm{Holm}}$ \\
\midrule
\endhead
All test & No change & 561 & -0.1704 & 1.1e-81 & 1.6e-80 \\
All test & Train mean & 561 & -0.0523 & 7.6e-47 & 9.9e-46 \\
All test & Additive & 561 & +0.0269 & 8.5e-23 & 9.4e-22 \\
All test & Linear (PCA) & 561 & -0.0222 & 3.2e-12 & 2.5e-11 \\
All test & Linear (ctrl emb.) & 561 & -0.0416 & 8.3e-46 & 9.9e-45 \\
All test & GEARS & 107 & +0.0167 & 0.0013 & 0.0025 \\
All test & GCP-Mamba ($\Delta$ only) & 561 & +0.0054 & 8.1e-07 & 4.8e-06 \\
All test & GCP-Mamba ($\Delta$ only, $\mathbf{w}_c\sim\mathcal{N}(0,1)$) & 561 & +0.0046 & 3.5e-05 & 0.00014 \\
All test & GCP-Mamba (permuted graph) & 561 & +0.0013 & 5.3e-05 & 0.00016 \\
All test & GCP-Mamba (random order) & 561 & +0.0064 & 1.5e-08 & 1e-07 \\
All test & Mamba (no graph) & 561 & +0.0024 & 4.1e-06 & 2.1e-05 \\
All test & Graph-MLP (no scan) & 561 & -0.0002 & 0.83 & 0.83 \\
All test & GCP-Mamba (additive prior) & 561 & +0.0273 & 5.2e-22 & 5.2e-21 \\
All test & Mamba (additive prior) & 561 & +0.0281 & 1.6e-21 & 1.4e-20 \\
\midrule
Unseen single & No change & 185 & -0.0540 & 1.7e-19 & 2.4e-18 \\
Unseen single & Train mean & 185 & -0.0023 & 0.06 & 0.42 \\
Unseen single & Additive & 185 & -0.0043 & 0.00087 & 0.0096 \\
Unseen single & Linear (PCA) & 185 & -0.0010 & 0.46 & 1 \\
Unseen single & Linear (ctrl emb.) & 185 & -0.0025 & 0.014 & 0.11 \\
Unseen single & GEARS & 36 & +0.0518 & 0.002 & 0.02 \\
Unseen single & GCP-Mamba ($\Delta$ only) & 185 & +0.0006 & 0.44 & 1 \\
Unseen single & GCP-Mamba ($\Delta$ only, $\mathbf{w}_c\sim\mathcal{N}(0,1)$) & 185 & -0.0003 & 0.35 & 1 \\
Unseen single & GCP-Mamba (permuted graph) & 185 & -0.0010 & 0.0088 & 0.079 \\
Unseen single & GCP-Mamba (random order) & 185 & +0.0014 & 0.15 & 0.74 \\
Unseen single & Mamba (no graph) & 185 & -0.0000 & 0.28 & 1 \\
Unseen single & Graph-MLP (no scan) & 185 & -0.0008 & 0.11 & 0.66 \\
Unseen single & GCP-Mamba (additive prior) & 185 & -0.0049 & 2.6e-05 & 0.00034 \\
Unseen single & Mamba (additive prior) & 185 & -0.0056 & 6.7e-05 & 0.0008 \\
\midrule
Double, 0/2 seen & No change & 32 & -0.2160 & 3.3e-08 & 4.6e-07 \\
Double, 0/2 seen & Train mean & 32 & -0.0487 & 0.0034 & 0.044 \\
Double, 0/2 seen & Additive & 32 & +0.0183 & 0.33 & 1 \\
Double, 0/2 seen & Linear (PCA) & 32 & -0.0417 & 0.018 & 0.2 \\
Double, 0/2 seen & Linear (ctrl emb.) & 32 & -0.0464 & 0.005 & 0.06 \\
Double, 0/2 seen & GEARS & 9 & +0.0941 & 0.055 & 0.55 \\
Double, 0/2 seen & GCP-Mamba ($\Delta$ only) & 32 & +0.0084 & 0.3 & 1 \\
Double, 0/2 seen & GCP-Mamba ($\Delta$ only, $\mathbf{w}_c\sim\mathcal{N}(0,1)$) & 32 & +0.0038 & 0.41 & 1 \\
Double, 0/2 seen & GCP-Mamba (permuted graph) & 32 & +0.0053 & 0.23 & 1 \\
Double, 0/2 seen & GCP-Mamba (random order) & 32 & +0.0058 & 0.61 & 1 \\
Double, 0/2 seen & Mamba (no graph) & 32 & -0.0019 & 0.75 & 1 \\
Double, 0/2 seen & Graph-MLP (no scan) & 32 & +0.0006 & 0.83 & 1 \\
Double, 0/2 seen & GCP-Mamba (additive prior) & 32 & +0.0172 & 0.32 & 1 \\
Double, 0/2 seen & Mamba (additive prior) & 32 & +0.0067 & 0.59 & 1 \\
\midrule
Double, 1/2 seen & No change & 254 & -0.2765 & 1.8e-40 & 2.5e-39 \\
Double, 1/2 seen & Train mean & 254 & -0.0995 & 3.9e-29 & 4.6e-28 \\
Double, 1/2 seen & Additive & 254 & +0.0575 & 2.2e-21 & 2.2e-20 \\
Double, 1/2 seen & Linear (PCA) & 254 & -0.0389 & 3.5e-08 & 2.8e-07 \\
Double, 1/2 seen & Linear (ctrl emb.) & 254 & -0.0800 & 8.8e-30 & 1.2e-28 \\
Double, 1/2 seen & GEARS & 43 & +0.0077 & 0.062 & 0.062 \\
Double, 1/2 seen & GCP-Mamba ($\Delta$ only) & 254 & +0.0106 & 0.00029 & 0.0011 \\
Double, 1/2 seen & GCP-Mamba ($\Delta$ only, $\mathbf{w}_c\sim\mathcal{N}(0,1)$) & 254 & +0.0086 & 0.00071 & 0.0021 \\
Double, 1/2 seen & GCP-Mamba (permuted graph) & 254 & +0.0046 & 2e-05 & 9.8e-05 \\
Double, 1/2 seen & GCP-Mamba (random order) & 254 & +0.0127 & 4e-06 & 2.8e-05 \\
Double, 1/2 seen & Mamba (no graph) & 254 & +0.0064 & 5.5e-06 & 3.3e-05 \\
Double, 1/2 seen & Graph-MLP (no scan) & 254 & +0.0020 & 0.03 & 0.06 \\
Double, 1/2 seen & GCP-Mamba (additive prior) & 254 & +0.0576 & 8.2e-22 & 9e-21 \\
Double, 1/2 seen & Mamba (additive prior) & 254 & +0.0610 & 8.1e-21 & 7.3e-20 \\
\midrule
Double, 2/2 seen & No change & 90 & -0.3615 & 1.7e-16 & 2.4e-15 \\
Double, 2/2 seen & Train mean & 90 & -0.2286 & 1.7e-16 & 2.4e-15 \\
Double, 2/2 seen & Additive & 90 & +0.1065 & 6.4e-15 & 7e-14 \\
Double, 2/2 seen & Linear (PCA) & 90 & -0.0720 & 2e-05 & 9.9e-05 \\
Double, 2/2 seen & Linear (ctrl emb.) & 90 & -0.1615 & 4.2e-15 & 5.1e-14 \\
Double, 2/2 seen & GEARS & 19 & +0.0015 & 0.57 & 0.57 \\
Double, 2/2 seen & GCP-Mamba ($\Delta$ only) & 90 & +0.0157 & 9.1e-05 & 0.00036 \\
Double, 2/2 seen & GCP-Mamba ($\Delta$ only, $\mathbf{w}_c\sim\mathcal{N}(0,1)$) & 90 & +0.0239 & 7e-06 & 4.9e-05 \\
Double, 2/2 seen & GCP-Mamba (permuted graph) & 90 & +0.0064 & 1.3e-05 & 7.5e-05 \\
Double, 2/2 seen & GCP-Mamba (random order) & 90 & +0.0150 & 0.00035 & 0.0011 \\
Double, 2/2 seen & Mamba (no graph) & 90 & +0.0175 & 5.6e-06 & 4.5e-05 \\
Double, 2/2 seen & Graph-MLP (no scan) & 90 & -0.0032 & 0.1 & 0.2 \\
Double, 2/2 seen & GCP-Mamba (additive prior) & 90 & +0.1045 & 7e-15 & 7e-14 \\
Double, 2/2 seen & Mamba (additive prior) & 90 & +0.1063 & 7.7e-15 & 7e-14 \\
\bottomrule
\end{longtable}


\begin{table}[h]
\caption{Variation due to random initialisation on the Norman data. For each split, the MSE (top-20 DE genes, all test perturbations) of three independently initialised runs; shown are the mean over splits of the within-split range and standard deviation, and the mean split-level MSE.}\label{tab:reps}
\centering\small
\begin{tabular}{lrrr}
\toprule
Model & Mean MSE & Within-split s.d. & Within-split range \\
\midrule
GCP-Mamba & 0.290 & 0.024 & 0.046 \\
GCP-Mamba (permuted graph) & 0.285 & 0.023 & 0.042 \\
Mamba (no graph) & 0.278 & 0.016 & 0.032 \\
Graph-MLP (no scan) & 0.288 & 0.017 & 0.033 \\
GCP-Mamba (additive prior) & 0.225 & 0.001 & 0.003 \\
\bottomrule
\end{tabular}
\end{table}


\begin{table*}[!t]
\processtable{Pearson correlation of predicted and observed expression change on the top-20 DE genes (mean\,$\pm$\,s.d.\ over five splits; higher is better)\label{tab:pearson}}{
\begin{tabular*}{\hsize}{@{\extracolsep{\fill}}lccccc@{}}
\toprule
Model & Unseen single & Double, 0/2 seen & Double, 1/2 seen & Double, 2/2 seen & All test \\
\midrule
No change & --- & --- & --- & --- & --- \\
Train mean & 0.53\,$\pm$\,0.05 & 0.72\,$\pm$\,0.12 & 0.66\,$\pm$\,0.06 & 0.67\,$\pm$\,0.07 & 0.62\,$\pm$\,0.04 \\
Additive & 0.54\,$\pm$\,0.05 & 0.70\,$\pm$\,0.15 & 0.82\,$\pm$\,0.03 & 0.97\,$\pm$\,0.01 & 0.75\,$\pm$\,0.03 \\
Linear (PCA) & 0.52\,$\pm$\,0.07 & 0.64\,$\pm$\,0.18 & 0.73\,$\pm$\,0.03 & 0.84\,$\pm$\,0.04 & 0.67\,$\pm$\,0.03 \\
Linear (ctrl emb.) & 0.52\,$\pm$\,0.05 & 0.70\,$\pm$\,0.09 & 0.69\,$\pm$\,0.06 & 0.76\,$\pm$\,0.08 & 0.65\,$\pm$\,0.05 \\
\midrule
GCP-Mamba & 0.55\,$\pm$\,0.08 & 0.71\,$\pm$\,0.13 & 0.76\,$\pm$\,0.03 & 0.86\,$\pm$\,0.04 & 0.71\,$\pm$\,0.04 \\
\midrule
GCP-Mamba ($\Delta$ only) & 0.57\,$\pm$\,0.06 & 0.71\,$\pm$\,0.15 & 0.76\,$\pm$\,0.02 & 0.87\,$\pm$\,0.03 & 0.72\,$\pm$\,0.03 \\
GCP-Mamba ($\Delta$ only, $\mathbf{w}_c\sim\mathcal{N}(0,1)$) & 0.56\,$\pm$\,0.07 & 0.72\,$\pm$\,0.14 & 0.76\,$\pm$\,0.03 & 0.88\,$\pm$\,0.03 & 0.71\,$\pm$\,0.04 \\
GCP-Mamba (permuted graph) & 0.53\,$\pm$\,0.08 & 0.68\,$\pm$\,0.16 & 0.76\,$\pm$\,0.03 & 0.86\,$\pm$\,0.03 & 0.70\,$\pm$\,0.04 \\
GCP-Mamba (random order) & \textbf{0.58\,$\pm$\,0.08} & \textbf{0.73\,$\pm$\,0.12} & 0.77\,$\pm$\,0.03 & 0.88\,$\pm$\,0.02 & 0.72\,$\pm$\,0.04 \\
Mamba (no graph) & 0.55\,$\pm$\,0.08 & 0.70\,$\pm$\,0.14 & 0.77\,$\pm$\,0.03 & 0.87\,$\pm$\,0.03 & 0.71\,$\pm$\,0.04 \\
Graph-MLP (no scan) & 0.55\,$\pm$\,0.06 & 0.68\,$\pm$\,0.17 & 0.77\,$\pm$\,0.02 & 0.86\,$\pm$\,0.03 & 0.70\,$\pm$\,0.03 \\
\midrule
GCP-Mamba (additive prior) & 0.53\,$\pm$\,0.06 & 0.70\,$\pm$\,0.17 & 0.82\,$\pm$\,0.03 & 0.98\,$\pm$\,0.01 & 0.74\,$\pm$\,0.03 \\
Mamba (additive prior) & 0.54\,$\pm$\,0.06 & 0.70\,$\pm$\,0.18 & \textbf{0.82\,$\pm$\,0.03} & \textbf{0.98\,$\pm$\,0.01} & \textbf{0.75\,$\pm$\,0.04} \\
\botrule
\end{tabular*}}{}
\end{table*}


\section{Hyperparameters and compute}
\begin{table}[h]
\caption{Hyperparameters shared by all deep-model variants.}\label{tab:hyper}
\centering\small
\begin{tabular}{ll}
\toprule
Parameter & Value\\
\midrule
Layers, model width $d$, state size $N$ & 2, 32, 4\\
Readout rank & 16\\
Gene embedding & 64 PCs of control-cell expression, standardised\\
Co-expression graph & $|r|$ on control cells, $k=20$ neighbours, symmetrised\\
Graph signal & 3 diffusion steps, decay 0.5, $\log(1+10^3\,\cdot)$\\
Scan chunk length, log-decay clamp & 16, $-5$\\
Optimiser & AdamW, lr $10^{-3}$, weight decay $10^{-2}$, batch 8, clip 1.0\\
Dropout (perturbation encoder) & 0.1\\
Epochs & min.\ 40, max.\ 100, patience 15 on validation MSE\\
Seeds & model seed = split seed (1--5)\\
Hardware & 4-core CPU, no GPU\\
\bottomrule
\end{tabular}
\end{table}


\end{document}
